% concatenated from bayesianWCS.tex
%%%%%%%%%%%%%%%%%%%%%%%%%%%%%%%%%%%%%%%%
%\documentclass[11pt,reqno,dvipdfmx]{amsart}

\documentclass[11pt,reqno,letter]{amsart}

\usepackage{amssymb}
%\usepackage{graphicx}
\usepackage[dvipdfmx]{graphicx}

\usepackage{placeins}

%\usepackage{hyperref}
%\usepackage{color}
%\usepackage{amsmath}
%\usepackage{bm}
\usepackage[usenames]{color}
\usepackage{caption}
%\usepackage{latexsym}
\usepackage[dvips]{psfrag}
%\usepackage{graphicx}
\usepackage{latexsym}
%\usepackage{amssymb}
\usepackage{amsmath}
\usepackage{amsthm}
\usepackage{tabularx}
%\usepackage{algorithm}
%\usepackage{algorithmic}
\usepackage{bm}
\usepackage{lscape}
\usepackage[numbers]{natbib}
\usepackage[usenames]{color}
\usepackage{url}
\usepackage{hyperref}
%\usepackage{textgreek}
\usepackage{upgreek}
\usepackage{units}
\usepackage{fancyhdr}
\definecolor{OliveGreen}{rgb}{0,0.6,0}

\usepackage{enumitem}

\usepackage[normalem]{ulem}

\renewcommand{\baselinestretch}{1.3}

\headheight=8pt \topmargin=0pt \textheight=624pt
\textwidth=432pt \oddsidemargin=18pt \evensidemargin=18pt

\setlength{\parskip}{\baselineskip}
\setlength{\parindent}{0pt}%

\theoremstyle{plain}
\newtheorem{theorem}{Theorem}[section]
\newtheorem{lemma}[theorem]{Lemma}
\newtheorem{constraint}[theorem]{Constraint}
\newtheorem{corollary}[theorem]{Corollary}
\newtheorem{proposition}[theorem]{Proposition}
\newtheorem{conjecture}[theorem]{Conjecture}
\newtheorem{result}[theorem]{Result}
\newtheorem{fact}{Fact}
\theoremstyle{definition}
\newtheorem{definition}[theorem]{Definition}
\newtheorem{remark}[theorem]{Remark}
\newtheorem{example}[theorem]{Example}
\newtheorem{assumption}[theorem]{Assumption}
\newtheorem{property}[theorem]{Property}

\numberwithin{equation}{section}
\numberwithin{theorem}{section}
\numberwithin{table}{section}
\numberwithin{figure}{section}

\DeclareMathOperator*{\argmax}{arg\,max}
\DeclareMathOperator*{\argmin}{arg\,min}
%%%%%%%%%%%%%%%%%%%%%%%%%%%%%%%%%%%%%%%%
\newcommand{\N}{{\rm I\!N}}
\newcommand{\R}{{\rm I\!R}}
\newcommand{\PI}{{\rm I\!\Pi}}
\newcommand{\e}{{\rm e}}
\newcommand{\E}{{\mathbb E}}
\newcommand{\p}{{\mathbb P}}
\newcommand{\V}{{\mathbb V}}
\newcommand{\q}{{\mathbb Q}}

%\let\footnote\thanks

%%%%%%%%%%%%%%%%%%%%%%%%%%%%%%%%%%%%%%%%%%%%%%%
%\title{Marginal Increase of Robustness of DRO Objectives\\ with Indifferentiable $\phi$-Divergences}
%\title{Worst-case sensitivity}
%\title[Mean-Sensitivity Tradeoffs in DRO]{Mean-Sensitivity Tradeoffs\\ in Distributionally Robust Optimization}

%\shorttitle{Mean-Sensitivity Tradeoffs in Distributionally Robust Optimization}
%\author{Gotoh, Kim, Lim}


%\title{Distributionally Robust Optimization is a Multi-Objective Problem}%\thanks{An earlier version of this paper went by the title ``Worst-Case Sensitivity".}
%\title[Robustness Measures in Distributionally Robust Optimization]{Robustness Measures in \\ Distributionally Robust Optimization}
%\title{Robustness measures for mixtures and Bayesian models}
%\thanks{An earlier version of this paper had the title ``Distributionally Robust Optimization is a Multi-Objective Problem"; prior to that it was ``Worst-Case Sensitivity".}

%\titlenote{An earlier version of this paper went by the title ``Distributionally Robust Optimization is a Multi-Objective Problem". An even earlier version went by ``Worst-Case Sensitivity". We are open to suggestions for future titles.}

\title[Posterior and Likelihood Sensitivity in Bayesian DRO]{Posterior and Likelihood Sensitivity in \\Bayesian Distributionally Robust Optimization}

\date{\today}

\author[Gotoh]{Jun-ya Gotoh\textsuperscript{$\dagger$}}

\author[Kim]{Michael Jong Kim\textsuperscript{$\ddagger$}}

\author[Lim]{Andrew E.B. Lim\textsuperscript{$*$}}

\dedicatory{\textsuperscript{$\dagger$}Department of Data Science for Business Innovation, Chuo University, Tokyo, Japan. Email: jgoto@kc.chuo-u.ac.jp \\ \textsuperscript{$\ddagger$}Sauder School of Business, University of British Columbia, Vancouver, Canada. Email: mike.kim@sauder.ubc.ca \\ \textsuperscript{$*$}Department of Analytics and Operations, Department of Finance, and Institute of Operations Research and Analytics, National University of Singapore, Singapore. Email: andrewlim@nus.edu.sg}
%%%%%%%%%%%%%%%%%%%%%%%%%%%%%%%%%%%%%%%%%%%%%%%

%%%%%%%%%%%%%%%%
\begin{document}
%%%%%%%%%%%%%%%%


\begin{abstract}
We introduce the notion of worst-case posterior and worst-case likelihood sensitivity. These measure, respectively, the sensitivity of the expected cost to worst-case perturbations of the posterior distribution and worst-case perturbations of the likelihood of a Bayesian model. Each defines a quantitative measure of robustness. A decision maker concerned about the sensitivity of the out-of-sample expected cost to deviations from her assumptions will want a decision for which both sensitivities are small. We derive posterior and likelihood sensitivities for uncertainty sets defined in terms of deviation measures. Posterior sensitivity vanishes when the posterior variance shrinks to zero, which occurs when parameter uncertainty is eliminated from learning. Parameter learning does not  eliminate likelihood sensitivity. A distributionally robust formulation of a Bayesian optimization problem makes a near-Pareto-optimal tradeoff between performance (expected cost) and robustness (posterior and likelihood sensitivity).



\end{abstract}


\maketitle
%%%%%%%%%%

{\bf Key words:} Bayesian distributionally robust optimization, worst-case sensitivity, model uncertainty, posterior sensitivity, likelihood sensitivity, robustness--performance tradeoff.


\section{Introduction}

%Worst--Case Sensitivity (WCS) is the sensitivity of the expected cost ${\mathbb E}_{\mathbb P}[f(x, Y)]$ under worst-case deviations from a nominal model $\mathbb P$ for the random variable $Y$. WCS is a quantitative measure of robustness: The expected cost is sensitive to misspecificaton if worst-case sensitivity is large. A decision maker (DM) concerned about robustness would like a decision $x$ with low sensitivity and (ideally) an expected cost that is not too much larger than the minimum. Though formulated as a worst-case problem, a large class of Distributionally Robust Optimization (DRO) problems map out a near-Pareto-optimal tradeoff between expected cost (performance) and WCS (robustness), where the sensitivity measure in the tradeoff is determined by the deviation measure ${\rm d}({\mathbb Q}|{\mathbb P})$ that defines the uncertainty set \cite{gotoh2026sensitivity}. This provides a basis for selecting the family and size of the uncertainty set according to the characteristics of the nominal model and the cost function and guidance for designing systems where the price of robustness is low.

Bayesian models provide a natural framework for learning from data. A Bayesian model requires two inputs from the decision maker (DM), a  likelihood function $L^\theta(\cdot) \equiv L(\cdot|\theta)$ and a prior distribution $\rho$ for the uncertain parameter $\theta$. Together they characterize the joint distribution of the uncertain parameter $\theta$ and the random variable $Y$. The prior on $\theta$ is updated using observed samples of $Y$. The resulting posterior   induces the predictive distribution which is used to compute the expected cost under a decision $x$. A DM  concerned about misspecification would like a small expected cost that is insensitive  to deviations from the predictive distribution.

In this paper, we introduce the notions of posterior and likelihood sensitivity. Posterior sensitivity is the sensitivity of the expected cost to worst-case deviations from the posterior while likelihood sensitivity defines an analogous notion for the likelihood function. They can be interpreted as quantitative measures of robustness.  A DM  concerned about the impact of misspecification on the out-of-sample performance  would like a decision with small expected cost and low posterior and likelihood sensitivity. 
%This paper builds  on the closely related idea of worst-case--sensitivity that has been used as a robustness measure for empirical optimization problems \cite{gotoh2018robust,gotoh2021calibration,gotoh2026sensitivity}. Here we adapt this idea  to the posterior and likelihood of a Bayesian problem.
Adapting ideas from \cite{gotoh2018robust,gotoh2021calibration,gotoh2026sensitivity} for robust empirical optimization problems, we show that a broad class of distributionally robust Bayesian optimization problems can be approximated by a ``regularized" Bayesian problem where the regularizer is a weighted sum of posterior and likelihood sensitivities. It follows that approximately Pareto optimal tradeoffs between performance (expected cost) and robustness (posterior and likelihood sensitivity) can be obtained by solving a Bayesian Distributionally Robust Optimization (DRO) problem. 
%We show that a broad class of distributionally robust Bayesian optimization problems can be approximated by a ``regularized" Bayesian problem where the regularizer is a weighted sum of posterior and likelihood sensitivities. It follows that a nearly Pareto-optimal tradeoff between performance (expected cost) and robustness (posterior and likelihood sensitivity) can be obtained by solving a Bayesian Distributionally Robust Optimization (DRO) problem. 

We derive expressions for posterior and likelihood sensitivity, which depend on the uncertainty sets for posterior and likelihood, and show that posterior sensitivity vanishes when the posterior distribution degenerates to a point mass, but not so for the likelihood. In other words, improvements in parameter estimates from updating the prior with data can reduce posterior, but not likelihood, sensitivity. As illustrated in \cite{gotoh2026sensitivity} for empirical optimization, sensitivity measures can guide the selection of uncertainty sets, identify when the ``price of robustness" is high, and guide system redesign to reduce this cost. 
% This means that posteriorfor example from updating the prior with a large amount of data. This is not the case for the likelihood function.

The outline of the paper is as follows. We briefly review the literature in Section \ref{sec:lit}. We introduce the nominal model (the predictive distribution resulting from the posterior and likelihood) in Section \ref{sec:nominal} and worst-case posterior and worst-case likelihood sensitivity in Section \ref{sec:WCS}. Examples where uncertainty sets for the posterior and likelihood are smooth $\phi$-divergence and formed using different $\phi$-divergences are considered. We show in Section \ref{sec:DROtradeoff} that Bayesian DRO problems map out a nearly Pareto-optimal tradeoff between expected cost, posterior sensitivity and likelihood sensitivity. We consider a simple pricing application in Section \ref{sec:experiments}. We adopt the ``constrained" formulation of the Bayesian DRO problem, where alternative posterior and likelihood distributions are determined by a constraint on the deviation from the nominal. A generalization to the penalty formulation of the Bayesian model can be found in the Appendix.

\section{Literature review}
\label{sec:lit}

This paper is related to three streams of work: approximations of DRO as a regularized nominal problem, robustness measures induced by worst-case sensitivity, and Bayesian DRO.

Several papers \cite{bertsimas2018,blanchet2019,Duchi2017variance,esfahani2018data,gao2023distributionally, gao2024wasserstein,kuhn2019wasserstein} show that DRO can be approximated by a regularized empirical optimization problem. \cite{gotoh2026sensitivity} takes this as the starting point, shows that the regularizer can be interpreted as a quantitative measure of robustness (worst-case sensitivity), and explores the implications for robust modeling including uncertainty set selection (family and size) and system design. The present paper extends this perspective to single stage Bayesian problems. 

A Bayesian model consists of a likelihood function and a prior distribution; the predictive distribution (formed from the posterior and likelihood) is the nominal model. Generalizing ideas from \cite{gotoh2026sensitivity}, we introduce worst-case posterior and worst-case likelihood sensitivity so that a Bayesian decision maker can quantify the fragility of her out-of-sample expected cost to the components of the predictive distribution. We show how each sensitivity measure depends on the uncertainty sets for the posterior and likelihood and that a distributionally robust Bayesian problem maps out a tradeoff between performance (expected cost) and robustness (posterior and likelihood sensitivities). Our results show that posterior sensitivity vanishes when there is no parameter uncertainty, but not so for the likelihood. \cite{shapiro2023bayesian} develops Bayesian DRO and, among other results, derive worst-case likelihood sensitivity  with a Kullback--Leibler uncertainty set. %It also establishes consistency of the posterior, and convergence of the objective and optimal Bayesian DRO solutions. 
We complement this by introducing posterior sensitivity, and considering uncertainty sets for deviation measures beyond Kullback--Leibler divergence.

There is a related body of work in Bayesian statistics that studies the sensitivity of posterior-based inferences  to perturbations of the (subjective) prior and  likelihood (see for example \cite{berger1986Bayesian,gustafson1985Bayesian,Lavine1991Bayesian}). A key difference is the focus on decision making and its role in our model (i.e., we explicitly include the objective function); in particular, how decisions can be chosen when posterior or likelihood sensitivity is large. By showing that the DRO objective can be decomposed into  expected cost plue a weighted sum of posterior sensitivity and likelihood sensitivity, we provide a systematic approach to choosing decisions that balance  performance and sensitivity to misspecification.


%\section{Mean--sensitivity Tradeoff}


%\section{Mixtures and Bayesian models}
%\label{App:Bayes}


%\section{Mixtures and Bayesian models} \label{sec:Bayesian}
\section{Nominal model}
\label{sec:nominal}

Let $(\theta, Y)$ be random variables with a joint distribution ${\mathbb P}({\rm d}\theta, {\rm d}y)$. We denote the marginal distribution ${\mathbb P}({\rm d}\theta) = \rho({\rm d}\theta)$ 
and the conditional distribution  ${\mathbb P}({\rm d}y|\theta) = L^\theta({\rm d}y)$ so
\begin{eqnarray*}
{\mathbb P}({\rm d}\theta, {\rm d}y) = \rho({\rm d}\theta)L^\theta({\rm d}y).
\end{eqnarray*}
Bayesian learning models belong to this class: ${\mathcal L} = \{L^{\theta}\vert \theta\in\Theta\}$ is a parameterized family of likelihoods and $\rho$ is the posterior distribution of the uncertain parameter $\theta$.
%For a Bayesian model, ${\mathcal L} = \{L^{\theta}(y)\vert \theta\in\Theta\}$ is a parameterized family of likelihood functions and  $\rho(d\theta)$ is the posterior distribution  of the uncertain parameter $\theta \in \Theta$, belong to this class. 
The posterior is obtained by updating prior $\rho_0(d\theta)$ using data $Y_1, \cdots, Y_m$ and Bayes' rule. The likelihood function $L^\theta$ is a modeling choice, while the posterior depends on the user-specified prior $\rho_0(d\theta)$ and likelihood $L^{\theta}(dy)$. We are concerned about the sensitivity of the expected cost to deviations from ${\mathbb P}({\rm d}\theta, {\rm d}y)$.  

Another important class is mixture models with populations indexed by $\theta\in\{\theta_1, \cdots, \theta_n\}$ and population distribution ${\mathbb P}({\rm d}y|\theta=\theta_i) = L^{\theta_i}({\rm d}y)$. 
For concreteness, we refer to the marginal distribution $\rho(d\theta)$ as the posterior and the conditional distribution $L^\theta({\rm d}y)$ as the likelihood, fully recognizing that it is more general than a Bayesian learning model. 

Let $f$ be a cost function of the variables $(\theta, Y)$. %For concreteness, we refer to $f(\theta, Y)$ as a cost. 
The expected cost is
\begin{eqnarray*}
{\mathbb E}_{\mathbb P}[f(\theta, Y)] = {\mathbb E}_\rho\big\{{\mathbb E}_{L^\theta}[f(\theta, Y)\, | \, \theta]\big\}.
\end{eqnarray*}
The expected cost depends on the joint distribution. We are concerned about the sensitivity of the expected cost relative to worst-case deviations from $\mathbb P$. 

%\subsection{Worst-case problems}

\section{Worst-case sensitivity} 
\label{sec:WCS}

We use ideas from \cite{gotoh2026sensitivity} to define worst-case sensitivity with respect to the posterior and the likelihood.  


\subsection{Uncertainty in the posterior and likelihood}

We begin by defining deviations from the posterior $\rho({\rm d}\theta)$ and likelihood ${\mathbb P}({\rm d}y | \theta)$. 

Let $\eta$ be a distribution  on $\Theta$  and ${\rm d}( \eta | \rho)$ be a deviation measure of $\eta$ from $\rho$. For every $\theta\in\Theta$ let ${\mathbb Q}^\theta$ be an alternative to $L^\theta$ for the distribution of  $Y$ and ${\rm d}({\mathbb Q}^\theta| L^\theta)$ a measure of deviation of ${\mathbb Q}^\theta$ from $L^\theta$. It will occasionally be convenient to write ${\mathcal L} = \{ L^\theta| \theta\in\Theta\}$ when we talk about the entire family of distributions.

%Let ${\mathcal Q} = \{{\mathbb Q}^\theta\vert\theta\in\Theta\}$ be a family of probability distributions for $Y$ parameterized by $\theta\in\Theta$ and ${\mathcal L} = \{ L^\theta| \theta\in\Theta\}$ be the nominal likelihood function. For every $\theta$, ${\rm d}({\mathbb Q}^\theta| L^\theta)$ is the deviation of ${\mathbb Q}^\theta$ from $L^\theta$. 
%We define the deviation of ${\mathcal Q}$ from $\mathcal L$ as the expected deviation of the probability measure ${\mathbb Q}^\theta$ from ${\mathbb P}^\theta$ under the posterior $\rho$:
%\begin{eqnarray}
%{\rm d}({\mathcal Q}| {\mathcal L}) & = & {\mathbb E}_\mathbb \rho \big\{{\rm d}({\mathbb Q}^\theta| L^\theta)\big\}.
%\label{eq:dQL}
%\end{eqnarray}



%\subsection{Constraint model}
Let
\begin{eqnarray}
 \label{eq:uncertainty-set}
{\mathcal P}(\varepsilon) & = & \big\{ \eta\big\vert {\rm d}(\eta|\rho)\leq \varepsilon\big\},\; \\  [5pt]
{\mathcal L}^\theta(\delta) & =& \big\{{\mathbb Q}\big\vert {\rm d}({\mathbb Q}|L^\theta)\leq\delta\big\} \nonumber
\end{eqnarray}
be uncertainty sets for the posterior and likelihood distributions. We use the same notation ${\rm d}(\cdot|\cdot)$ for the deviation measure in each uncertainty set. However, different deviation measures can be used for the posterior and likelihood (see Section \ref{sec:mix}).

Consider the worst-case expected cost
\begin{eqnarray}
V(\epsilon, \delta) 
%& = & \min_{\eta \in {\mathcal P}_\phi(\varepsilon)} \min_{\{{\mathbb Q}^\theta\} \in {\mathcal L}_\phi(\delta)} {\mathbb E}_\eta {\mathbb E}_{Q^\theta} [f(x, Y)] \\
= \max_{\eta \in {\mathcal P}(\varepsilon)} {\mathbb E}_\eta \Big\{\max_{{\mathbb Q}^\theta \in {\mathcal L}^\theta(\delta)} {\mathbb E}_{Q^\theta} [f(\theta, Y)] \Big\}.
\label{eq:wc_Bayes}
\end{eqnarray}
The inner expectation is conditional on $\theta$; we consider worst-case deviations ${\mathbb Q}^\theta$ from the nominal likelihood $L^\theta$ for each $\theta$. The outer expectation considers worst-case deviations $\eta$ from the posterior $\rho$ so changes the weights on $\theta$. 

\subsection{Regularization}

Consider first the inner problem. Since the worst-case expected cost is increasing in the size of the uncertainty set $\delta$, we can write for every $\theta\in\Theta$
\begin{eqnarray}
\Psi(\theta) = \max_{{\mathbb Q}^\theta \in {\mathcal L}^\theta(\delta)} {\mathbb E}_{{\mathbb Q}^\theta} [f(\theta, Y)] 
   =    {\mathbb E}_{L^\theta}[f(\theta, Y)]  + {\mathcal A}_{L^\theta}\big(\delta; f(\theta, Y)\big)
   \label{eq:inner}
   \end{eqnarray}
where   ${\mathcal A}_{L^\theta}\big(\delta; f(\theta, Y)\big)$, the ambiguity cost, is non-negative and non-decreasing in $\delta$. 

Similarly, we can write the outer problem
\begin{eqnarray}
\max_{\eta \in {\mathcal P}(\varepsilon)} {\mathbb E}_\eta \big\{\Psi(\theta) \big\} & = & {\mathbb E}_\rho \big\{\Psi(\theta) \big\} + {\mathcal A}_{\rho}\big(\varepsilon; \Psi(\theta)\big)
\label{eq:outer}
\end{eqnarray}
where the ambiguity cost ${\mathcal A}_{\rho}\big(\varepsilon; \Psi(\theta)\big)$ is again non-negative and increasing in $\varepsilon$.


Substituting \eqref{eq:inner} into \eqref{eq:outer} it  follows that the worst-case expected cost \ref{eq:wc_Bayes} can be written
\begin{eqnarray*}
V(\epsilon, \delta)  = {\mathbb E}_\rho \Big\{{\mathbb E}_{L^\theta}[f(\theta, Y)]  \Big\} + {\mathcal A}_{\rho,{\mathcal L}}\big(\varepsilon, \delta; f(\theta, Y)\big)
\end{eqnarray*}
with ambiguity cost
\begin{eqnarray*}
{\mathcal A}_{\rho,{\mathcal L}}\big(\varepsilon, \delta; f(\theta, Y)\big) := {\mathbb E}_\rho \Big\{ {\mathcal A}_{L^\theta}\big(\delta; f(\theta, Y)\big) \Big\} + {\mathcal A}_{\rho}\Big(\varepsilon; {\mathbb E}_{L^\theta}[f(\theta, Y)]  + {\mathcal A}_{L^\theta}\big(\delta; f(\theta, Y)\big)\Big).
\end{eqnarray*}



We make the following assumptions about the ambiguity costs. The functions $g(\varepsilon)$ and $g(\delta)$ for the two ambiguity costs need not be the same.
\begin{assumption} \label{ass:Bayes}
The ambiguity costs ${\mathcal A}_{L^\theta}\big(\delta; f(\theta, Y)\big)$  and ${\mathcal A}_{\rho}\big(\varepsilon; \Psi(\theta)\big)$ are such that 
\begin{enumerate}
\item there is a non-decreasing function $g(\delta)$ such that $g(\delta)\rightarrow 0$ when $\delta\rightarrow 0$ and 
\begin{eqnarray*}
%\Psi^\delta(\theta, f(\theta, \cdot)) & =& 
%\max_{{\mathbb Q}^\theta \in {\mathcal L}(\delta)} {\mathbb E}_{Q^\theta} [f(\theta, Y)] 
 %  & = &   {\mathbb E}_{L^\theta}[f(\theta, Y)]  + 
{\mathcal A}_{L^\theta}\big(\delta; f(\theta, Y)\big) %\nonumber \\ [5pt]
   =  g(\delta) {\mathcal S}_{L^\theta}(f(\theta, Y))+ o(g(\delta))
%\label{eq:Psi}
\end{eqnarray*}
for every $\theta\in{\Theta}$;
%where 
%\begin{eqnarray*}
\item there is a non-decreasing function $g(\varepsilon)$ such that $g(\varepsilon)\rightarrow 0$ when $\varepsilon\rightarrow 0$ and 
\begin{eqnarray*}
%\max_{\eta \in {\mathcal P}(\varepsilon)} {\mathbb E}_\eta \big\{\Psi(\theta) \big\} & = & {\mathbb E}_\rho \big\{\Psi(\theta) \big\} + 
{\mathcal A}_{\rho}\big(\varepsilon; \Psi(\theta)\big) %\\ [5pt]
=  %{\mathbb E}_\rho \big\{\Psi(\theta) \big\} + 
g(\varepsilon) {\mathcal S}_\rho(\Psi(\theta)) + o(g(\varepsilon)).
\end{eqnarray*}
\item  ${\mathcal S}_{\rho}(\Psi)$ is continuous in $\Psi$. 
\end{enumerate}
\end{assumption}
Intuitively, $\Psi(\theta)$ is the worst-case expected cost when  the uncertainty set  is $\{{\mathcal L}^\theta(\delta)\}$ and
\begin{eqnarray*}
{\mathcal A}_{L^\theta}\big(\delta; f(\theta, Y)\big)= \max_{{\mathbb Q}^\theta \in {\mathcal L}^\theta(\delta)} {\mathbb E}_{{\mathbb Q}^\theta} [f(\theta, Y)] 
   -   {\mathbb E}_{L^\theta}[f(\theta, Y)]
   \end{eqnarray*} 
is the  increase in expected cost relative to the nominal. When Assumption \ref{ass:Bayes} holds,  ${\mathcal S}_{L^\theta}(f(\theta, Y))$  can be interpreted as the  sensitivity of the expected cost under worst-case perturbations; it depends on the deviation measure ${\rm d}({\mathbb Q} | L^\theta)$ that defines the uncertainty set. For example, ${\mathcal S}_{L^\theta}(f(\theta, Y))$ is the standard deviation of the cost $\sigma_{L^\theta} (f(\theta, Y))$ when ${\rm d}({\mathbb Q}|L^\theta)$ is smooth $\phi$-divergence. Sensitivity measures induced by other measures of deviation can be found in \cite{gotoh2026sensitivity} and references cited therein (see also Section \ref{sec:mix}).
   
Likewise, ${\mathbb E}_\rho \big\{\Psi(\theta) \big\}$ is the expected value of $\Psi(\theta)$ when $\theta$ has distribution $\rho$ and
${\mathcal S}_{\rho}(\Psi)$ is the sensitivity of the expected cost ${\mathbb E}_\rho \big\{\Psi(\theta) \big\}$ under worst-case perturbations of $\rho$. 
   
Uncertainty in posterior and likelihood contribute to increases in the total expected cost
   \begin{eqnarray*}
   {\mathcal A}_{\rho,{\mathcal L}}\big(\varepsilon, \delta; f(\theta, Y)\big) = \max_{\eta \in {\mathcal P}(\varepsilon)} {\mathbb E}_\eta \Big\{\max_{{\mathbb Q}^\theta \in {\mathcal L}^\theta(\delta)} {\mathbb E}_{Q^\theta} [f(\theta, Y)] \Big\} - {\mathbb E}_\rho \Big\{{\mathbb E}_{L^\theta}[f(\theta, Y)]  \Big\}.
   \end{eqnarray*}
Our goal is to understand how the sensitivity of the total cost depends on the posterior and likelihood.

The following result gives an expansion of the value function when $\varepsilon$ and $\delta$ are small. 
%The proof can be found in Appendix \ref{App:Bayes}.
\begin{theorem} \label{prop:Bayesian expansion}
Suppose that Assumption \ref{ass:Bayes} holds. Then
\begin{eqnarray*}
V(\epsilon, \delta) & = & 
{\mathbb E}_\rho\Big\{{\mathbb E}_{L^\theta}[f(\theta, Y)] \Big\} + g(\delta){\mathbb E}_\rho\Big\{{\mathcal S}_{L^\theta}(f(\theta, Y))\Big\}  +  g(\varepsilon) {\mathcal S}_\rho\Big({\mathbb E}_{L^\theta}(f(\theta, Y))\Big) \\ [5pt] 
& & + o(g(\delta)) + o(g({\varepsilon})).
\end{eqnarray*}
\end{theorem}


%\subsection{Proof of Proposition \ref{prop:Bayesian expansion}}
\begin{proof}
It follows from \eqref{eq:inner} and Assumption \ref{ass:Bayes} that
\begin{eqnarray*}
\Psi(\theta)  &= &   {\mathbb E}_{L^\theta}[f(\theta, Y)]  + {\mathcal A}_{L^\theta}\big(\delta; f(\theta, Y)\big) \\ & = & {\mathbb E}_{L^\theta}[f(\theta, Y)]  + g(\delta) {\mathcal S}_{L^\theta}(f(\theta, Y))+ o(g(\delta)).
\end{eqnarray*}
From \eqref{eq:outer} and Assumption \ref{ass:Bayes} we have
\begin{align*}
V(\varepsilon, \delta) & =  \max_{\eta \in {\mathcal P}(\varepsilon)} {\mathbb E}_\eta \big\{\Psi(\theta) \big\} \\[8pt] 
& = \max_{\eta \in {\mathcal P}(\varepsilon)} {\mathbb E}_\eta\Big\{{\mathbb E}_{L^\theta}[f(\theta, Y)]  + g(\delta) {\mathcal S}_{L^\theta}(f(\theta, Y))+ o(g(\delta))\Big\} \\[8pt]
& = {\mathbb E}_\rho\Big\{{\mathbb E}_{L^\theta}[f(\theta, Y)]+ g(\delta) {\mathcal S}_{L^\theta}(f(\theta, Y))+ o(g(\delta)) \Big\} \\[8pt]
& \qquad + g(\varepsilon){\mathcal S}_\rho\Big({\mathbb E}_{L^\theta}[f(\theta, Y)]  + g(\delta) {\mathcal S}_{L^\theta}(f(\theta, Y))+ o(g(\delta))\Big) + o(g(\varepsilon)) \\[8pt]
& = {\mathbb E}_\rho\Big\{{\mathbb E}_{L^\theta}[f(\theta, Y)] \Big\} + g(\delta){\mathbb E}_\rho\Big\{{\mathcal S}_{L^\theta}(f(\theta, Y))\Big\}  +  g(\varepsilon) {\mathcal S}_\rho\Big({\mathbb E}_{L^\theta}(f(\theta, Y))\Big) \\[8pt]
& \qquad + g(\varepsilon) \left\{{\mathcal S}_\rho\Big({\mathbb E}_{L^\theta}[f(\theta, Y)]  + g(\delta) {\mathcal S}_{L^\theta}(f(\theta, Y)) + o(g(\delta)) \Big) - {\mathcal S}_\rho\Big({\mathbb E}_{L^\theta}(f(\theta, Y))\Big)\right\}  
\\[8pt]
& 
\qquad + o(g(\delta)) + o(g({\varepsilon})) \\[8pt]
& = {\mathbb E}_\rho\Big\{{\mathbb E}_{L^\theta}[f(\theta, Y)] \Big\} + g(\delta){\mathbb E}_\rho\Big\{{\mathcal S}_{L^\theta}(f(\theta, Y))\Big\}  +  g(\varepsilon) {\mathcal S}_\rho\Big({\mathbb E}_{L^\theta}(f(\theta, Y))\Big) \\[8pt] 
& \qquad + o(g(\delta)) + o(g({\varepsilon}))
\end{align*}
where the last equality follows from the continuity of ${\mathcal S}_\rho(\Psi)$.
\end{proof}
%\hfill$\Box$



\subsection{Worst-case sensitivity}

Given likelihood function $L^\theta(y)$ ($\theta\in\Theta$) and posterior $\rho$, 
worst-case posterior sensitivity is the sensitivity of the expected cost under worst-case deviations from the posterior defined by the uncertainty set ${\mathcal P}(\varepsilon)$:
\begin{eqnarray*}
{\mathcal S}_\rho(f)  =  \lim_{\epsilon\rightarrow 0}\frac{V(\epsilon, 0)- {\mathbb E}_\rho\Big\{{\mathbb E}_{L^\theta}[f(\theta, Y)]\Big\}}{g({\epsilon})}% \nonumber \\ [5pt]
=
{\mathcal S}_\rho\Big\{{\mathbb E}_{L^\theta}[f(\theta, Y)]\Big\}.
%\label{eq:prior-sensitivity-const}
\end{eqnarray*}

Given the nominal posterior $\rho$, likelihood  $L^\theta(y)$, and uncertainty set ${\mathcal L}^\theta(\delta)$ for deviations from the likelihood, worst-case likelihood sensitivity is the sensitivity of the expected cost with respect to worst-case deviations from the nominal likelihood
\begin{eqnarray*}
{\mathcal S}_{\mathcal L}(f)  =  \lim_{\delta\rightarrow 0}\frac{V(0, \delta)- {\mathbb E}_\rho\Big\{{\mathbb E}_{L^\theta}[f(\theta, Y)]\Big\}}{g(\delta)}% \nonumber \\ [5pt]
 =  {\mathbb E}_\rho \Big\{{\mathcal S}_{L^\theta}(f(\theta, Y))\Big\}.
%\label{eq:likelihood-sensitivity-const}
\end{eqnarray*}
Note that posterior (likelihood) sensitivity ${\mathcal S}_\rho$  depends on the deviation measure ${\rm d}(\eta|\rho)$  that defines the uncertainty set ${\mathcal P}(\varepsilon)$, and likelihood sensitivity  ${\mathcal S}_{L^\theta}$ on the deviation measure  ${\rm d}({\mathbb Q}|L^\theta)$ that controls deviations from the nominal likelihood. 





In general, ${\mathcal S}_\rho$ is a generalized measure of deviation \cite{gotoh2026sensitivity,rockafellar2006generalized} so posterior sensitivity is zero when $\rho$ is a point mass (i.e., there is no uncertainty about $\theta$) or when $g(\theta) = {\mathbb E}_{L^\theta}[f(\theta, Y)]$ is constant on the support of $\rho$. %Likelihood sensitivity is non-zero unless $f(\theta, y)$ is constant as a function of $y$ on the support of $L^\theta$. 
Likelihood sensitivity is zero only when, for $\rho$-almost every $\theta$, the sensitivity of the cost under $L^\theta$ is zero. For the deviation measures considered here, this occurs when $f(\theta, y)$ is constant in $y$ on the support of $L^\theta$.



The following example considers posterior and likelihood sensitivity when deviation measures are smooth $\phi$-divergence.

\subsection{Smooth $\phi$-divergence}
 \label{eg:phi_Bayes}
 

Suppose that $\eta$ is absolutely continuous with respect to $\rho$ and  ${\mathbb Q}^\theta$ is absolutely continuous with respect to $L^\theta$ for every $\theta\in\Theta$. Let $\frac{{\rm d} \eta}{{\rm d}\rho}$ denote the Radon--Nikodym derivative of $\eta$ with respect to $\rho$, and for every $\theta\in\Theta$, $\frac{{\rm d}Q^\theta}{{\rm d}L^\theta}$ be the Radon--Nikodym derivative of $Q^{\theta}$ with respect to $L^\theta$.


For the likelihood function,  $\phi$-divergence of ${\mathbb Q}^\theta$ with respect to $L^\theta$, for every $\theta\in{\Theta}$, is defined by
\begin{eqnarray*}
 {\rm d}({\mathbb Q}^\theta|L^\theta)= {\mathbb E}_{L^\theta} \left[\phi \left(\frac{{\rm d}{\mathbb Q}^\theta}{{\rm d} L^\theta}\right) \right].
\end{eqnarray*}
In the case of smooth $\phi$-divergence, Assumption \ref{ass:Bayes} holds and the inner problem \cite{gotoh2026sensitivity} 
\begin{eqnarray*}
\max_{{\mathbb Q}^\theta \in {\mathcal L}^\theta_{\phi}(\delta)} {\mathbb E}_{Q^\theta} [f(\theta, Y)] 
=  {\mathbb E}_{L^\theta}[f(\theta, Y)]  + \sqrt{\frac{2\delta}{\phi''(1)}} {\sigma}_{L^\theta}(f(\theta, Y))+ o(\sqrt{\delta}).
\end{eqnarray*}
For the outer problem we measure the deviation of $\eta$ from $\rho$ using $\phi$-divergence:
\begin{eqnarray*}
{\rm d}( \eta | \rho) & = & {\mathbb E}_\rho \left[\phi\left(\frac{{\rm d}\eta}{{\rm d}\rho}\right)\right]
\end{eqnarray*}
and it follows that
\begin{eqnarray*}
\max_{\eta \in {\mathcal P}_\phi(\varepsilon)} {\mathbb E}_\eta[\Psi(\theta)] 
= {\mathbb E}_\rho[\Psi(\theta)] + \sqrt{\frac{2\varepsilon}{\phi''(1)}}\sigma_\rho(\Psi(\theta)) + o(\sqrt{\varepsilon}).
\end{eqnarray*}
In particular, $g(\delta) = \sqrt{\delta}$ and 
\begin{eqnarray*} 
{\mathcal S}_{L^\theta}(f) =  \sqrt{\frac{2}{\phi''(1)}}  {\sigma}_{L^\theta}(f(\theta, Y))
\end{eqnarray*}  
for the inner problem, where ${\sigma}_{L^\theta}(f(\theta, Y)) = \sqrt{ {\mathbb V}_{L^\theta}(f(\theta, Y))}$ is the standard deviation of the cost under the distribution $L^\theta$. For the outer problem,
$g(\varepsilon) = \sqrt{\varepsilon}$ and
\begin{eqnarray*}
{\mathcal S}_\rho(\Psi(\theta))  =\sqrt{\frac{2}{\phi''(1)}}\sigma_\rho(\Psi(\theta)).
\end{eqnarray*}
Note that the regularizer in the expansion of the worst-case expected costs are given by the standard deviation because we are using smooth $\phi$-divergence as the deviation measure. Other uncertainty sets give different regularizers \cite{gotoh2026sensitivity}.


It now  follows from Theorem \ref{prop:Bayesian expansion} that the worst-case expected cost 
\begin{eqnarray*}
V(\epsilon, \delta) %& = & \min_{\eta \in {\mathcal P}_\phi(\varepsilon)} {\mathbb E}_\eta \Big\{\min_{{\mathbb Q}^\theta \in {\mathcal L}_\phi(\delta)} {\mathbb E}_{Q^\theta} [f(x, Y)] \Big\}  \label{eq:WC-Bayes-penalty} \\
& = &  {\mathbb E}_\rho\Big\{{\mathbb E}_{L^\theta}[f(\theta, Y)]\Big\}+ \sqrt{\frac{2\delta}{\phi''(1)}}{\mathbb E}_\rho \Big({\sigma}_{L^\theta}(f(\theta, Y))\Big) + \sqrt{\frac{2\varepsilon}{\phi''(1)}}
{\sigma}_\rho\Big({\mathbb E}_{L^\theta}[f(\theta, Y)]\Big) \nonumber \\
& & + o(\sqrt{\epsilon}) + o(\sqrt{\delta}).
%\label{eq:WC-Bayes-penalty}
\end{eqnarray*}
Worst-case posterior  sensitivity is
\begin{eqnarray}
{\mathcal S}_\rho(f) %& = & \lim_{\epsilon\rightarrow 0}\frac{V(\epsilon, 0)- {\mathbb E}_\rho\Big\{{\mathbb E}_{L^\theta}[f(x, Y)]\Big\}}{\sqrt{\epsilon}} \nonumber \\ 
%&  = & 
%\sqrt{\frac{2}{\phi''(1)}}
= \sqrt{\frac{2}{\phi''(1)}}{\sigma}_\rho\Big({\mathbb E}_{L^\theta}[f(\theta, Y)]\Big)
\label{eq:prior-sensitivity-const}
\end{eqnarray}
and worst-case likelihood sensitivity is %worst-case sensitivity with respect to the likelihood is defined by
\begin{eqnarray}
{\mathcal S}_{\mathcal L}(f) %& = & \lim_{\epsilon\rightarrow 0}\frac{V(0, \delta)- {\mathbb E}_\rho\Big\{{\mathbb E}_{L^\theta}[f(x, Y)]\Big\}}{\sqrt{\delta}} \nonumber \\ 
%&  = &
= \sqrt{\frac{2}{\phi''(1)}} {\mathbb E}_\rho \Big\{{\sigma}_{L^\theta}(f(\theta, Y))\Big\}.
\label{eq:likelihood-sensitivity-const}
\end{eqnarray}
Both sensitivity measures depend on the posterior and objective function, but quite differently. Posterior sensitivity  is the standard deviation of the conditional expectation $g(\theta) = {\mathbb E}_{L^\theta}[f(\theta, Y)]$ and will be small if $g(\theta)$ is ``flat" on the support of $\rho$, even if the standard deviation of $\theta$ is large. 
It vanishes when the posterior variance shrinks (e.g.) because of learning. Likelihood sensitivity depends on the variance of the cost given $\theta$. It typically will not vanish even if $\theta$ is known because there is still uncertainty about the distribution $L^\theta$ of $Y$.


The paper \cite{shapiro2023bayesian} derives worst-case likelihood sensitivity  with a Kullback--Leibler uncertainty set (the inner ``max" in \eqref{eq:wc_Bayes}) but not the posterior. We complement this by introducing posterior sensitivity and by showing how the worst-case objective can be approximated by  expected cost, posterior sensitivity, and likelihood sensitivity.


\subsection{Mix and match}
\label{sec:mix}

As noted in the discussion following \eqref{eq:uncertainty-set}, there is no reason to restrict ourselves to smooth $\phi$-divergence for the deviation measure or to have the same deviation measure for the posterior and likelihood. For example, if instead of smooth $\phi$-divergence, we use 
\begin{eqnarray*}
\phi(z)=\delta_{[1-\varepsilon,\frac{1-(1-\varepsilon)\alpha}{1-\alpha}]}(z)
\end{eqnarray*} 
to define the uncertainty set for the likelihood function, we have \cite{gotoh2026sensitivity}\footnote{Given $\theta$, $\mathrm{CVaR}_{L^\theta,\alpha}(f(\theta, Y))$ is the conditional-value-at-risk of the cost $f(\theta, Y)$ at level $\alpha$ when $Y$ has distribution $L^\theta$: \begin{eqnarray*}\mathrm{CVaR}_{L^\theta,\alpha}(f(\theta, Y)):=\min_\gamma\Big\{\gamma+\frac{1}{1-\alpha}\mathbb{E}_{L^\theta}\Big(\max\big\{f(\theta, Y)-\gamma,0\big\}\Big)\Big\}\end{eqnarray*}}
\begin{eqnarray*} 
{\mathcal S}_{L^\theta}(f) = \mathrm{CVaR}_{L^\theta,\alpha}(f(\theta, Y))-\mathbb{E}_{L^\theta} (f(\theta, Y)).
\end{eqnarray*} 
(WCS for other deviation measures are also derived in \cite{gotoh2026sensitivity}). We choose this uncertainty set and deviation measure if we are concerned about misspecification of the right tail of the cost distribution. 
If smooth $\phi$-divergence is used for the posterior and this alternative is used for the likelihood, we have
\begin{align*}
{\mathcal S}_{\rho}(f) &= \sqrt{\frac{2}{\phi''(1)}}{\sigma}_\rho\Big({\mathbb E}_{L^\theta}[f(\theta, Y)]\Big), \\[5pt]
{\mathcal S}_{\mathcal L}(f) & = {\mathbb E}_\rho \Big\{\mathrm{CVaR}_{L^\theta, \alpha}(f(\theta, Y))-\mathbb{E}_{L^\theta} (f(\theta, Y))\Big\}.
\end{align*}



%Posterior sensitivity is important because it quantifies the sensitivity of the expected cost to worst-case deviations from the posterior which depends on the user-specified prior distribution. 

%\subsection{Other uncertainty sets}

%An important class of deviation measures are of the form
%\begin{eqnarray*}
%{\rm d}({\mathbb Q}|{\mathbb P}) =  {\mathbb E}_{\mathbb P}\Big[\phi\Big(\frac{d {\mathbb Q}}{d {\mathbb P}}\Big)\Big].
%\end{eqnarray*}
%Another important deviation measure is the Wasserstein distance. By selecting ${\mathbb P} \equiv \rho$ or ${\mathbb P} \equiv L^\theta$ these deviation measures can be used to form uncertainty sets posterior and likelihood sensitivity. 
%Worst-case sensitivity for different choices of $\phi$ are provided in Table \ref{table:summary} as well as for the Wasserstein metric. Although smooth $\phi$-divergence  is used for both the posterior and likelihood uncertainty sets in Section \ref{eg:phi_Bayes}, other deviation measures can be used and it is not necessary for the deviation measure for the posterior uncertainty set to be the same as that of the likelihood (as noted in the discussion immediately following \eqref{eq:uncertainty-set}).


%\begin{table}
%\begin{center}
%\caption{Worst-case sensitivities considered in Section \ref{sec:WCS}}
%\label{table:summary}
%\begin{tabular}{ll|c}
%\multicolumn{3}{c}{$\langle$i$\rangle$~ $\max_{\mathsf{q}}\{\mathbb{E}_{\mathsf{q}}(\mathsf{f})\vert d(\mathsf{q}\vert\mathsf{p})\leq \varepsilon\}$}\\
%\hline
%\multicolumn{2}{l|}{Type of divergence/worst-case objective} & Worst-case sensitivity \eqref{eq:sensitivity-general2} \\
%\hline
%\multicolumn{2}{l|}{(a) smooth $\phi$-divergence
%\qquad
%$\phi''(1)>0=\phi(1)=\phi'(1)$} & $\sqrt{\frac{2{\mathbb V}_{\mathsf{p}}(\mathsf{f})}{\phi''(1)}}$ \\ [2pt]
%(b) total variation &
%$\phi(z)=|z-1|$ & $\frac{1}{2}\big(\max(\mathsf{f})-\min(\mathsf{f})\big)$\\ [2pt]
%(c) ``budgeted'' 
%& $\phi=\delta_{[0,1+\varepsilon]}$ & $\mathbb{E}_{\mathsf{p}}(\mathsf{f})-\min(\mathsf{f})$ \\[2pt]
%(d) $(1-\varepsilon)$``mean% cost
%''$+\varepsilon$``$\alpha$-CVaR'' & 
%$\phi=\delta_{[1-\varepsilon,\frac{1-(1-\varepsilon)\alpha}{1-\alpha}]}$
%& $\mathrm{CVaR}_{\mathsf{p},
%\alpha}(\mathsf{f})-\mathbb{E}_{\mathsf{p}}
%(\mathsf{f})$\\ [5pt]
%(d') $(1-\varepsilon)$``mean
%''$+\varepsilon$``max
%'' 
%& $\phi=\delta_{[1-\varepsilon,
%\infty)}$
% & $\max(\mathsf{f})-\mathbb{E}_{\mathsf{p}}(\mathsf{f})$\\[2pt]
%(d'') symmetric ($\alpha=\frac{1}{2-\varepsilon}$) & 
%$\phi=\delta_{[1-\varepsilon,\frac{1}{1-\varepsilon}]}$ & $\mathrm{CVaR}_{\mathsf{p},\frac{1}{2}}(\mathsf{f})-\mathbb{E}_{\mathsf{p}}(\mathsf{f})$ \\[5pt]
%\multicolumn{2}{l|}{(e) type-1 Wasserstein distance 
%~w/ Lipschitz smooth $f$ 
%} & $\max\limits_{i=1,\cdots,n} \max\limits_{z_i}\frac{f(z_i)-f(Y_i)}{\|z_i-Y_i\|}$ \\
%\hline
%\end{tabular}\\
%\end{center}
%\end{table}



\section{Distributionally Robust Optimization}
\label{sec:DRO}

\subsection{Performance--sensitivity tradeoffs in DRO}
\label{sec:DROtradeoff}

Theorem \ref{prop:Bayesian expansion} shows that the  worst-case expected cost can be locally expanded as
\begin{eqnarray}
\lefteqn{\min_x \max_{\eta \in {\mathcal P}(\varepsilon)} {\mathbb E}_\eta \Big\{\max_{{\mathbb Q}^\theta \in {\mathcal L}^\theta(\delta)} {\mathbb E}_{Q^\theta} [f(x, \theta, Y)] \Big\}}
\label{eq:DRO-regularized}
\\ [5pt]
& = & \min_x {\mathbb E}_\rho\Big\{{\mathbb E}_{L^\theta}[f(x, \theta, Y)]\Big\}+ g(\delta) {\mathbb E}_\rho \Big({\mathcal S}_{L^\theta}(f(x, \theta, Y))\Big) + g(\varepsilon)  {\mathcal S}_\rho\Big({\mathbb E}_{L^\theta}[f(x, \theta, Y)]\Big) 
\nonumber  \\ [5pt] 
& & + o(g({\epsilon})) + o(g({\delta})).
\nonumber
\end{eqnarray}
This representation is related to results showing the relationship between DRO and a regularized nominal problem. For empirical optimization problems, \cite{gotoh2026sensitivity} shows that the regularization term is not just a mathematical object that connects the nominal and worst-case problems, but  has the physical interpretation as worst-case sensitivity. \eqref{eq:DRO-regularized} shows an analogous result for a Bayesian DRO problem, where the regularizer is now a weighted sum of posterior  sensitivity ${\mathcal S}_\rho({\mathbb E}_{L^\theta}[f(x, \theta, Y)])$   and likelihood sensitivity ${\mathbb E}_\rho ({\mathcal S}_{L^\theta}(f(x, \theta, Y)))$, and varying the uncertainty set sizes maps out a tradeoff between performance (expected cost) and robustness (posterior and likelihood sensitivity). 

%It is now interesting to understand how a robust DM makes tradeoffs between expected cost and posterior and likelihood sensitivity.





%\begin{remark}
%Worst-case prior and likelihood sensitivities for the penalty-version of the Bayesian DRO problem with a smooth $\phi$-divergence penalty is given in Appendix \ref{App:BayesPenalty}.
%\end{remark}







\section{Experiment}
\label{sec:experiments}


Demand $D(p)$ is normal with standard deviation $\sigma$ and expected value ${\mathbb E}[D(p)] = a + b p$ which depends on the price $p$. We assume that $\sigma$ is known and constant whereas $a$ and $b$ are constant but unknown to the DM. At the time of the decision, the decision maker has a posterior on $(a, b)$ which we assume to be normal. (This could have been obtained by updating an initial prior using data and Bayes' rule.) The DM's revenue when price is $p$ and observed demand is $D(p)$ is $p\cdot D(p)$. The expected revenue under the posterior is ${\mathbb E}_\rho [p \, D(p)] = p \cdot ({\mathbb E}_\rho[a]  + {\mathbb E}_\rho[b]\cdot p)$; the nominal DM chooses price to maximize this quantity. 

We assume a modified $\chi^2$--uncertainty set for the posterior and likelihood, i.e., $\phi(z) = \frac{1}{2}(z-1)^2$. It follows that posterior and likelihood sensitivities are given by
\begin{align}
\label{eq:pricing-senitivities}
{\mathcal S}_\rho[p\, D(p)] & = \sqrt{2} p \, \sigma_\rho [a + b p], \\
{\mathcal S}_{\mathcal L}[p\, D(p)] & = \sqrt{2}\sigma p. \nonumber
\end{align}
For this experiment, under the baseline posterior, $a$ has mean $5$ and standard deviation $0.75$, and $b$ has mean $-0.5$ and standard deviation $2.5$; we assume $a$ and $b$ are independent under the posterior. We set $\sigma = 3$.  For the nominal problem the optimal price is $p =\$5$; the optimal expected reward is $\$12.50$.

Figure \ref{fig:frontier} shows reward--posterior sensitivity and reward--likelihood sensitivity frontiers when the posterior variance for $a$ and $b$ is 0, $1/4$ of the baseline, the baseline, and double the baseline. This simulates the effect of learning, which reduces the posterior variance, on posterior and likelihood sensitivity. We see in the first plot that posterior sensitivity is increasing in the posterior variance, quantifying the sensitivity of the expected reward to the posterior. Note that posterior sensitivity is $0$ when there is no parameter uncertainty. The DRO solution makes a tradeoff between expected reward and worst-case posterior sensitivity. For this particular example, DRO is quite effective when the posterior variance is large ($2\times$baseline) as substantial reduction in the posterior sensitivity from its value under the nominal model can be achieved with minimal impact on expected cost. It is less effective at reducing posterior sensitivity when the posterior variance is small ($0.25 \times$baseline). However, the frontier shows that this is not really necessary as the posterior sensitivity is already relatively small. %We move along the frontier by changing the size of the uncertainty set; the associated DRO prices are shown in Figure \ref{fig:DRO_price}.

\begin{figure}[!htbp]
\begin{center}
\includegraphics[scale=0.26]{prior_sensitivity_vary_prior_var.png}  \\ [5pt]
\includegraphics[scale=0.26]{likelihood_sensitivity_vary_prior_var.png} 
%\includegraphics[scale=0.07]{frontier1.png} 
\end{center}
\caption{Reward--posterior sensitivity frontiers (upper), and reward--likelihood sensitivity frontiers (lower), for several posterior variances. Uncertainty sets for posterior and likelihood uncertainty are equal in size (i.e., $\delta = \varepsilon$); points  corresponding to $\varepsilon=0, 0.25$ and $0.5$ are indicated in the top plot. The expected reward under the nominal ($\varepsilon=\delta=0$) is $12.5$. From  the upper plot, posterior--sensitivity of the nominal solution is increasing in the  posterior variance; it vanishes when posterior variance is $0$. For this example, likelihood sensitivity is independent of the posterior variance.  }
\label{fig:frontier}
\end{figure}





From the lower plot, we see that the mean--likelihood sensitivity frontiers lie on top of one another (the small deviations are from numerical errors). This is not a general property but a quirk of this example; it would not be the case if $\sigma$ was also uncertain\footnote{
It is easy to see that
\begin{eqnarray*}
{\mathbb E}_\rho[p\,D(p)] = \frac{{\mathcal S}_{\mathcal L}}{\sigma\sqrt{2}}{\mathbb E}_\rho[a] + \left(\frac{{\mathcal S}_{\mathcal L}}{\sigma\sqrt{2}}\right)^2 {\mathbb E}_\rho[b] 
\end{eqnarray*}
so the expected reward, given likelihood sensitivity ${\mathcal S}_{\mathcal L}$, does not depend on the posterior variance. It follows that the reward--likelihood sensitivity frontiers do not depend on the posterior variance. It can also be shown that the reward--posterior sensitivity frontier does not depend on the likelihood variance $\sigma^2$ (again, a quirk of the problem).}. The main message, however, is that we  tradeoff between expected reward, posterior sensitivity, and likelihood sensitivity as the size of the uncertainty set $\varepsilon=\delta$ changes, with DRO prices being mapped  in Figure  \ref{fig:DRO_price}.

\begin{figure}[!htbp]
\begin{center}
\includegraphics[scale=0.26]{DRO_price.png}  
\end{center}
\caption{Optimal robust price as a function of the size of the uncertainty set ($\varepsilon=\delta$) for different posterior variance.
}
\label{fig:DRO_price}
\end{figure}

We can see from \eqref{eq:pricing-senitivities} that the likelihood sensitivity does not vanish when the posterior variance is $0$, so long as $p\neq 0$. This is a general property: while more data shrinks the posterior and eliminates  posterior sensitivity (as observed in Figure \ref{fig:frontier}),  uncertainty about the distribution of demand associated with the likelihood model remains unless $\sigma=0$.
 

 \begin{figure}[!htbp]
\begin{center}
\includegraphics[scale=0.26]{likelihoodvar-reward-prior-frontier.png} 
\\ [5pt]
\includegraphics[scale=0.25]{likelihoodvar-reward-likelihood-frontier.png} 
%\includegraphics[scale=0.07]{frontier1.png} 
\end{center}
\caption{ The upper figure shows expected reward versus posterior sensitivity whereas the lower figure shows expected reward versus likelihood sensitivity for a collection of likelihood variances $\sigma^2$. Uncertainty sets for posterior and likelihood are equal in size (i.e., $\delta = \varepsilon$). In this example, the reward–posterior sensitivity frontier is independent of the likelihood variance, while the reward–likelihood sensitivity frontier changes with $\sigma^2$.
%In contrast to Figure \ref{fig:frontier}, the reward--posterior sensitivity frontier is independent of the likelihood variance. 
}
\label{fig:likelihoodvar-frontier}
\end{figure}
 

Figure \ref{fig:likelihoodvar-frontier} shows performance--robustness frontiers when the posterior variance is held fixed and the variance $\sigma^2$ of the likelihood function is varied. The reward--likelihood sensitivity frontier is now increasing in $\sigma^2$ whereas the reward--posterior frontier remains unchanged.





%\begin{figure}[h]
%\begin{center}
%\includegraphics[scale=0.35]{likelihood_sen_eps.png} 
%\end{center}
%\caption{Likelihood sensitivity, for different prior variances, as a function of $\varepsilon$.}
%\label{fig:likelihood_sen}
%\end{figure}



%\clearpage

%\FloatBarrier

\section{Conclusion}
\label{sec:conclusion}

Bayesian models account for parameter uncertainty by specifying a prior on the uncertain parameters. However, the expected cost can still be sensitive to deviations from the  likelihood and the posterior (obtained by updating the user-specified prior using Bayes' rule and data) which may result in poor out-of-sample performance. Worst-case posterior sensitivity and worst-case likelihood sensitivity measure the sensitivity of the expected cost to worst-case deviations from the posterior and the likelihood, respectively. They are measures of robustness and a decision maker, concerned about the fragility of her assumptions, may prefer a decision for which both robustness measures are small. Worst-case posterior sensitivity vanishes when the posterior concentrates to a single point, which can occur when parameter uncertainty is reduced from learning. This is not the case with likelihood sensitivity. This is clear in the case of smooth $\phi$-divergence, where posterior sensitivity is the posterior standard deviation of the conditional expectation of the cost whereas likelihood sensitivity is the posterior expected value of its conditional standard deviation.
A nearly Pareto-optimal tradeoff between performance (expected cost) and robustness (worst-case posterior and likelihood sensitivity) can be mapped out by solving a Distributionally Robust Bayesian Optimization problem. 

%\clearpage



\paragraph{\bf Acknowledgments}
Jun-ya Gotoh is supported by the MEXT Grant-in-Aid 24K01113.
Michael Kim is supported  by the Natural Sciences and Engineering Research Council (NSERC) Discovery Grant RGPIN-2015-04019.
Andrew Lim is supported by the  Ministry of Education, Singapore, under its  2024 Academic Research Fund Tier 2 grant call (Award ref: MOE-T2EP20224-0018).


\newpage

\bibliographystyle{plainnat}
\bibliography{references}

\newpage


\appendix


\section{Penalty version of the Bayesian model}
\label{App:BayesPenalty}
To compute worst-case sensitivity for the penalty version in the case of smooth $\phi$-divergence, we proceed as before by defining worst-case deviations from the nominal distribution by solving a worst-case problem, and computing the sensitivity using the worst-case distribution. We account for the Bayesian structure by considering perturbations of the posterior and likelihood, which allows us to quantify the impact of these choices on sensitivity. 

Let $\eta$ be a distribution for $\theta$ on $\Theta$  and ${\rm d}( \eta | \rho)$ be a deviation measure of $\eta$ from $\rho$. 


Let ${\mathcal L} = \{ L^\theta| \theta\in\Theta\}$ be the nominal likelihood function and ${\mathcal Q} = \{{\mathbb Q}^\theta\vert\theta\in\Theta\}$ be a family of probability distributions for $Y$ parameterized by $\theta\in\Theta$. For every $\theta$, ${\rm d}({\mathbb Q}^\theta| L^\theta)$ is the deviation of ${\mathbb Q}^\theta$ from $L^\theta$. 
We define the deviation of ${\mathcal Q}$ from $\mathcal L$ as the expected deviation of the probability measure ${\mathbb Q}^\theta$ from ${\mathbb P}^\theta$ under the posterior $\rho$:
\begin{eqnarray}
{\rm d}({\mathcal Q}| {\mathcal L}) =  {\mathbb E}_\mathbb \rho \big\{{\rm d}({\mathbb Q}^\theta| L^\theta)\big\}.
\label{eq:dQL}
\end{eqnarray}



We find worst-case perturbations from the nominal model $\rho$ and $L^\theta$ by solving the worst-case problem
\begin{align*}
\lefteqn{\max_{{\eta, \,{\mathcal Q}} } {\mathbb E}_\eta \Big[{\mathbb E}_{{\mathbb Q}^{\theta}} \big\{f(x, Y)\big\}\Big] - \frac{1}{\delta}{\rm d}({\mathcal Q}| {\mathcal L}) - \frac{1}{\varepsilon} {\rm d}(\eta | \rho)} \\
& = \max_\eta  {\mathbb E}_\eta\left[\max_{{{\mathbb Q}^\theta\in {\mathcal Q}} }{\mathbb E}_{{\mathbb Q}^{\theta}} \big\{f(x, Y)\big\} - \frac{1}{\delta}{\mathbb E}_{L^\theta}\left\{\phi\left(\frac{{\rm d}Q^\theta}{{\rm d}L^\theta}\right)\right\}\right] - \frac{1}{\varepsilon}{\mathbb E}_\rho \left[\phi\left(\frac{{\rm d}\eta}{{\rm d}\rho}\right)\right]  \\
& = \max_\eta  {\mathbb E}_\eta \big[\Psi(\theta)\big] - \frac{1}{\varepsilon}{\mathbb E}_\rho \left[\phi\left(\frac{{\rm d}\eta}{{\rm d}\rho}\right)\right],
\end{align*}
where
\begin{eqnarray*}
\Psi(\theta) =\frac{1}{\delta} \ln {\mathbb E}_{L^{\theta}} \Big[e^{\delta f(x, Y)}\Big]
\end{eqnarray*}
and $\varepsilon$ controls deviations of the alternative posterior $\eta$ from the nominal $\rho$, while $\delta>0$ controls deviations of ${\mathbb Q}^\theta$ from $L^\theta$.
The worst-case distribution is given by
\begin{eqnarray*}
{\mathbb Q}^{\varepsilon,\delta}({\rm d}\theta, {\rm d}y) = \eta^\varepsilon({\rm d}\theta) {\mathbb Q}^{\theta,\delta}({\rm d}y),
\end{eqnarray*}
where
\begin{align*}
{\mathbb Q}^{\theta, \delta}  & =   \argmax_{\mathbb Q}{\mathbb E}_{{\mathbb Q}} \big[f(x, Y) \big] - \frac{1}{\delta}{\mathbb E}_{L^\theta}\left[\phi\left(\frac{{\rm d}{\mathbb Q}}{{\rm d}L^\theta}\right)\right],\\
\eta^\varepsilon & =  \argmax_\eta  {\mathbb E}_\eta \big[\Psi(\theta)\big] - \frac{1}{\varepsilon}{\mathbb E}_\rho \left[\phi\left(\frac{{\rm d}\eta}{{\rm d}\rho}\right)\right].
\end{align*}
The following result characterizes the worst-case distribution when $\delta$ and $\eta$ are small.
\begin{proposition} \label{prop:wc-bayes}
The worst-case measure ${\mathbb Q}^{\varepsilon, \delta}({\rm d}\theta, {\rm d}y)$
%\begin{eqnarray*}
%{\mathbb Q}^{\epsilon, \delta}({\rm d}\theta, dy) = \eta^\varepsilon({\rm d}\theta){\mathbb Q}^{\theta, \delta}(dy)
%\end{eqnarray*}
where
\begin{align*}
\frac{{\rm d}{\mathbb Q}^{\theta, \delta}}{{\rm d}L^{\theta}}  
& = [\phi']^{-1}\Big(\delta\big( f(Y)  -{\mathbb E}_{L^{\theta}}[f(Y)]\big)\Big) \\ 
& = 1 + \frac{\delta}{\phi{''}(1)}\Big\{f(Y)-{\mathbb E}_{L^\theta}[f(Y)]\Big\} + o(\delta),\\
\frac{{\rm d}\eta^\varepsilon}{{\rm d}\rho}
& = [\phi']^{-1}\Big(\varepsilon\big( \Psi(\theta)  -{\mathbb E}_{\rho}[\Psi(\theta)]\big)\Big) \\
& = 1+ \frac{\varepsilon}{\phi''(1)}\Big(\Psi(\theta)-{\mathbb E}_\rho[\Psi(\theta)]\Big).
\end{align*}
\end{proposition} 
We can now compute worst-case sensitivity. The expected reward under the worst-case measure is
\begin{align*}
V(\varepsilon, \delta; f(x,\cdot))  = {\mathbb E}_{{\mathbb Q}^{\varepsilon, \delta}}\big[f(x, Y)\big] 
  = {\mathbb E}_{\eta^{\varepsilon}} \Big[{\mathbb E}_{{\mathbb Q}^{\theta,\delta}} \big\{f(x, Y)\big\}\Big].
\end{align*}
By Proposition \ref{prop:wc-bayes}, the inner expectation
\begin{align*}
{\mathbb E}_{{\mathbb Q}^{\theta,\delta}} \big[f(x, Y)\big] & = {\mathbb E}_{L^\theta}\Big[\frac{{\rm d}{\mathbb Q}^{\theta, \delta}}{{\rm d}L^{\theta}} f(x, Y)\Big] \\
& = {\mathbb E}_{L^\theta}\Big[\Big(1 + \frac{\delta}{\phi{''}(1)}\Big\{f(Y)-{\mathbb E}_{L^\theta}[f(Y)]\Big\} \Big)f(x, Y)\Big]+o(\delta) \\
& = {\mathbb E}_{L^\theta}[f(x, Y)] + \frac{\delta}{\phi{''}(1)} {\mathbb V}_{L^{\theta}}[f(x, Y)] + o(\delta).
\end{align*}
It now follows that 
\begin{align*}
\lefteqn{V(\varepsilon, \delta; f(x, \cdot))}\\ 
%& =  &{\mathbb E}_{\eta^{\varepsilon}} \Big[{\mathbb E}_{{\mathbb Q}^{\theta,\delta}} \big\{f(x, Y)\big\}\Big] \\  
%V(\varepsilon, \delta; f(x, \cdot))
& = {\mathbb E}_\rho\Big[\frac{{\rm d}\eta^{\varepsilon}}{{\rm d}\rho}\Big({\mathbb E}_{L^\theta}[f(x, Y)] + \frac{\delta}{\phi{''}(1)} {\mathbb V}_{L^{\theta}}[f(x, Y)] \Big)\Big] + o(\delta) \\
& = {\mathbb E}_\rho \Big[{\mathbb E}_{L^\theta}\big\{f(x, Y)\big\}\Big]+ \frac{\delta}{\phi{''}(1)}  {\mathbb E}_{\rho} \Big[{\mathbb V}_{L^\theta}\big\{f(x, Y)\big\}\Big] + \frac{\varepsilon}{\phi{''}(1)} {\mathbb V}_{\rho} \Big[{\mathbb E}_{L^\theta}\big\{f(x, Y)\big\}\Big]  + O(\varepsilon \delta),
\end{align*}
and the worst-case sensitivity with respect to the posterior and the likelihood are given by
\begin{align}
{\mathcal S}_{\rho}(f) & = \lim_{\varepsilon\rightarrow 0}\frac{1}{\varepsilon} \left\{V(\varepsilon, 0; f(x, \cdot)) - {\mathbb E}_\rho \Big[{\mathbb E}_{L^\theta}\big\{f(x, Y)\big\}\Big]\right\}  \nonumber \\ 
& =  \frac{1}{\phi{''}(1)} {\mathbb V}_{\rho} \Big[{\mathbb E}_{L^\theta}\big\{f(x, Y)\big\}\Big]
\label{eq:WCSprior} 
\end{align}
and
\begin{align}
{\mathcal S}_{\mathcal L}(f) 
& = \lim_{\delta\rightarrow 0} \frac{1}{\delta} \left\{V(0, \delta; f(x, \cdot)) - {\mathbb E}_\rho \Big[{\mathbb E}_{L^\theta}\big\{f(x, Y)\big\}\Big]\right\}  \nonumber \\ 
& = \frac{1}{\phi{''}(1)} {\mathbb E}_{\rho} \Big[{\mathbb V}_{L^\theta}\big\{f(x, Y)\big\}\Big].
\label{eq:WCSlikelihood}
\end{align}


\end{document}

%%%%%%%%%%%%%%%%%%%%%%%%%%%%%%%%%%%%%%%%%%%%%%%%

\section{Worst-case sensitivity}
\label{sec:DRO_sensitivity}

%We introduce worst-case sensitivity (WCS) as a measure of robustness and discuss some of its properties. We propose a new formulation of robust decision making as a tradeoff between expected cost and WCS (robustness). Though intuitively appealing, mean--sensitivity objectives can be difficult to optimize due to nonlinearities and non-convexities introduced by the sensitivity measure. We  show, however, that (near) Pareto-optimal solutions can be computed by solving a worst-case DRO problem. WCS is defined in terms of an uncertainty set. In subsequent sections, we derive explicit expressions for WCS for commonly used uncertainty sets. This enables us to select uncertainty sets  for the DRO problem that reflect the sensitivity control we desire. The size of the uncertainty set can be chosen to achieve the desired mean--sensitivity tradeoff.

Worst-case sensitivity can be interpreted as a measure of robustness \cite{gotoh2026sensitivity}. Interpreted as such, DRO can be interpreted as a tradeoff between expected cost (performance) and worst-case sensitivity (robustness).

%\subsection{Worst-case sensitivity}


\subsection*{Nominal model.} Consider a cost function $f(x, Y)$ where $x$ is the decision and $Y$ is a  random variable with (nominal) distribution ${\mathbb P}$.  We would like to find a decision that performs well out-of-sample. 
A natural candidate is the solution of the nominal/Sample Average Approximation (SAA) problem
\begin{align*}
\min_x ~ {\mathbb E}_{\mathbb P}[f(x, Y)].
%\label{eq:pop_opt}
\end{align*}
However, out-of-sample performance of this decision can be poor if the  real-world distribution of $Y$ is different from the nominal $\mathbb P$ and the expected reward under the nominal is sensitive to these errors. 




\subsection*{Robustness measure: Worst-case sensitivity} 


%In order to  tradeoff between expected cost and worst-case sensitivity, worst-case sensitivity needs to be defined. 
Let ${\mathcal Q}(\varepsilon)$ be an {\it uncertainty set} of size $\varepsilon$, specifically,    a set of probability distributions increasing in $\varepsilon$ that contains the nominal distribution for every $\varepsilon\geq 0$ and degenerates to the nominal when $\varepsilon=0$. The worst-case expected cost with respect to ${\mathcal Q}(\varepsilon)$ is
\begin{align}
V(\varepsilon; f(x,\cdot)) &:=
\max_{{\mathbb Q}\in{\mathcal Q}(\varepsilon)} \mathbb{E}_{\mathbb{Q}}[f(x, Y)].
\label{eq:V}
\end{align}
To ease the notation, we often write $V(\varepsilon, x) \equiv V(\varepsilon; f(x,\cdot))$.  

Suppose  there is a non-negative, concave and increasing function\footnote{The function $g(\varepsilon)$ depends on the uncertainty set. $V(\varepsilon; f(x, \cdot))$, and hence the difference ${\mathcal A}\big(\varepsilon; f(x, \cdot)\big) :=V(\varepsilon, x) - {\mathbb E}_{\mathbb P}[f(x, Y)]$, is concave in $\varepsilon$ if the set of alternative probability measures can be written in the form ${\mathcal Q}(\varepsilon) = \{{\mathbb Q}\,|\, {\rm d}({\mathbb Q}|{\mathbb P})\leq \varepsilon\}$ for some convex function $d$ of ${\mathbb Q}$, which is the case for almost all models in the literature.  $V(\varepsilon)-V(0)$ is $O(\sqrt \varepsilon)$  and $g(\varepsilon) = \sqrt{\varepsilon}$ for smooth $\phi$-divergence and $O(\varepsilon$) for budgeted uncertainty so $g(\varepsilon)=\varepsilon$. } $g(\varepsilon)$ such that 
\begin{eqnarray*}
V(\varepsilon, x) - {\mathbb E}_{\mathbb P}[f(x, Y)]  \sim O(g(\varepsilon))
\end{eqnarray*}
for every fixed  $x$.
Worst-case sensitivity 
\begin{eqnarray}
{\mathcal S}_{\mathbb P}[f(x, \cdot)] =  \begin{displaystyle} \lim_{\varepsilon\downarrow 0}\frac{V(\varepsilon, \, x)-{\mathbb E}_{\mathbb P}[f(x, Y)]}{g(\varepsilon)} \end{displaystyle}
\label{eq:sensitivity-general2}
\end{eqnarray}
measures the rate the expected cost increases under small perturbations of the nominal model and hence can be interpreted as a measure of robustness. 
WCS is determined by the  uncertainty set ${\mathcal Q}(\varepsilon)$ so selecting an uncertainty set is like selecting a robustness measure. Sensitivity measures for a selection of uncertainty sets are derived in \cite{gotoh2026sensitivity}.




%%%%%%%

%\subsubsection*{A non-worst-case alternative} A decision is robust to model misspecification if its expected reward is insensitive to worst-case deviations from the nominal model. Taking this at face value,  we consider a multi-objective problem where the decision maker trades off between expected cost and worst-case sensitivity defined in \eqref{eq:sensitivity-general2}
%\begin{eqnarray}
 %\min_x {\mathbb E}_{\mathbb P} [f(x,\,Y)] + g(\varepsilon)\,{\mathcal S}_{\mathbb P}[f(x,\cdot)].
 %\label{eq:WCSmult}
%\end{eqnarray}
%The intuitive appeal of this approach is that it reflects the tradeoff between performance and robustness  of concern to the decision maker. A potential downside is that worst-case sensitivity may be a difficult to optimize. Returning to Figure \ref{fig:illus}, every decision $x$ maps to a point $(m, {\mathcal S}_{\mathbb P})$ in  mean--sensitivity space illustrated in the plot in the right. The solutions of the multi-objective problem \eqref{eq:WCSmult} maps out the mean--sensitivity frontier illustrated in the figure on the right.

%%%%%%%


\subsection*{Distributionally Robust Optimization (DRO)} DRO accounts for misspecification in $\mathbb P$ by optimizing worst-case expected cost over a set of probability measures ${\mathcal Q}(\varepsilon)$ that includes $\mathbb P$
\begin{align}
\min_x V(\varepsilon; f(x,\cdot)) &\equiv \min_x \max_{{\mathbb Q}\in{\mathcal Q}(\varepsilon)} \mathbb{E}_{\mathbb{Q}}[f(x, Y)].
\label{eq:DRO}
\end{align}
%One advantage of DRO is computational tractability; it does not, however, have the intuitive appeal of the multi-objective problem \eqref{eq:WCSmult}. 

%It turns out, however, that we can get the best of both worlds because solutions of the single-objective DRO problem \eqref{eq:DRO} are (almost) Pareto optimal for the multi-objective problem \eqref{eq:WCSmult}. In other words, the DRO problem enables us to explore the mean--sensitivity tradeoff in \eqref{eq:WCSmult} in a computationally tractable way. 

Since   ${\mathcal Q}(\varepsilon)$ is increasing in $\varepsilon$ and contain the nominal $\mathbb P$, the worst-case expected cost \eqref{eq:V} is monotonically increasing in $\varepsilon$ so
\begin{eqnarray*}
V(\varepsilon, x) = {\mathbb E}_{\mathbb P}[f(x, Y)] + {\mathcal A}\big(\varepsilon; f(x, \cdot)\big)
%\label{eq:DRO mean ambiguity-cost}
\end{eqnarray*}
where the {\it ambiguity cost}  ${\mathcal A}\big(\varepsilon; f(x, \cdot)\big)$ is a non-negative increasing function of $\varepsilon$ such that ${\mathcal A}\big(0; f(x, \cdot)\big)=0$. 
%(Figure \ref{fig:sensitivity-plots1}) 
For conveniently chosen uncertainty sets,  ${\mathcal A}\big(\varepsilon; f(x, \cdot)\big)$ can be characterized through the dual problem of \eqref{eq:V}.  
By \eqref{eq:sensitivity-general2} we have 
\begin{eqnarray*}
{\mathcal A}\big(\varepsilon; f(x, \cdot)\big) =V(\varepsilon, x)  -  {\mathbb E}_{\mathbb P}[f(x, Y)] =  g(\varepsilon)\,{\mathcal S}_{\mathbb P}[f(x,\cdot)] + o(g(\varepsilon))
\end{eqnarray*}
so the worst-case problem \eqref{eq:DRO}  is locally equivalent to a tradeoff between expected cost and worst-case sensitivity:
%induces a tradeoff between expected cost and worst-case sensitivity %can be written as is almost the same as the multi-objective problem \eqref{eq:WCSmult} with expected cost and worst-case sensitivity as objectives
\begin{eqnarray}
\label{eq:DRO-mean-sensitivity}
\min_x \max_{{\mathbb Q}\in{\mathcal Q}(\varepsilon)} \mathbb{E}_{\mathbb{Q}}[f(x, Y)] =  \min_x {\mathbb E}_{\mathbb P} [f(x,\,Y)] + g(\varepsilon)\,{\mathcal S}_{\mathbb P}[f(x,\cdot)] + o(g(\varepsilon)), \; \varepsilon>0.
\end{eqnarray}
%{\color{blue}While this extends existing ``regularized--nominal" approximations to a broader class of uncertainty sets, our main contribution is not this approximation per se, but the modeling consequences of interpreting the induced regularizer as a robustness measure.}


\subsection*{Sensitivity is a measure of deviation}

%The objectives in a multi-objective problem  reflect the tradeoffs in the problem. For example, the choice of risk measure in a mean--risk problem determine properties of the solution. For this reason, it is natural to understand the properties of worst-case sensitivity. 

Since DRO makes a tradeoff between performance and robustness \eqref{eq:DRO-mean-sensitivity}, properties of the robustness measure (WCS) determine those of the DRO solution. WCS for different uncertainty sets are derived in \cite{gotoh2026sensitivity}. We summarize general properties.



We begin with the  definition of a {\it Generalized Measure of Deviation} \cite{rockafellar2006generalized}. This  generalizes the notion of standard deviation and measures the spread of a random variable \cite{rockafellar2006generalized}. %and \ref{sec:Bayesian}.
\begin{definition}\label{def:deviation}
Let $f$ be a  random variable.
 ${\mathcal H}[f]$ is a {\it Generalized Measure of Deviation} or a {\it Generalized Measure of Spread} of $f$ if
\begin{enumerate}
\item ${\mathcal H}[f] \geq 0$ with equality if and only if $f$ is constant;
\item ${\mathcal H}[\beta f] = \beta{\mathcal H}[f]$ for every constant $\beta\geq 0$;
\item ${\mathcal H}[\alpha + f]={\mathcal H}[f]$ for every constant $\alpha\in{\mathbb R}$.
\end{enumerate}
\end{definition}



The following result shows that worst-case sensitivity ${\mathcal S}_{\mathbb P}[f(x, \cdot)]$ is a {\it generalized measure of deviation}, and hence, measures the spread of the cost distribution under the nominal.  
\begin{proposition}
\label{prop:deviation}
Let ${f}:\Omega\rightarrow{\mathbb R}$ be a random variable on $(\Omega, {\mathcal F})$ and let ${\mathbb P}$ a probability measure on this space. Suppose that ${\rm d}({\mathbb Q}\,|\,{\mathbb P})$ is convex and continuous over the set of probability measures $\mathbb Q$ which are absolutely continuous with respect to $\mathbb P$ and that ${\rm d}({\mathbb Q}|{\mathbb P})=0$ if and only if ${\mathbb Q}={\mathbb P}$. 
Suppose too that there is a constant $k>0$ such that for every probability measure ${\mathbb Q}^{(\delta)}$ such that
\begin{eqnarray*}
\frac{\rm{d} {\mathbb Q}}{\rm{d} {\mathbb P}}^{(\delta)} = 1 + \delta {\Delta} > 0
\end{eqnarray*}
for some random variable $\Delta: \Omega \rightarrow {\mathbb R}$ satisfying ${\mathbb E}_{\mathbb P}(\Delta)=0$, that
\begin{align}
{\rm d}({\mathbb Q}^{(\delta)}\,|\,{\mathbb P}) & \sim O(\delta^k) ~\mbox{when}  ~\delta \rightarrow 0.
\label{eq:continuity d}
\end{align}
Let ${\mathcal A}(\varepsilon; f) $ be the ambiguity cost \eqref{eq:AC} for an uncertainty set
%\begin{eqnarray*}
%{\mathcal A}(\varepsilon; f) & = & \max_{{\mathbb Q}\in{\mathcal Q}(\varepsilon)}{\mathbb E}_{\mathbb Q} \big[f(Y)\big] - {\mathbb E}_{\mathbb P}\big[f(Y)\big]
%\end{eqnarray*}
%be the ambiguity cost where
\begin{eqnarray*}
{\mathcal Q}(\varepsilon)    =  \Big\{{\mathbb Q}:{\mathcal F} \rightarrow {\mathbb R} ~\Big|~{\rm d}({\mathbb Q}\,| \,{\mathbb P})\leq\varepsilon\Big\}.
\end{eqnarray*}
%is the uncertainty set.
 For every $\varepsilon>0$, ${\mathcal A}(\varepsilon;f)$ is a  generalized measure of deviation of $f$,
%If  there is a constant $k>0$ such that
%\begin{align}
%{\rm d}({\mathbb Q}^{(\delta)}\,|\,{\mathbb P}) & \sim O(\delta^k) ~\mbox{when}  ~\delta \rightarrow 0,
%\label{eq:continuity d}
%\end{align}
%for every probability measure ${\mathbb Q}^{(\delta)}$ such that
%\begin{eqnarray*}
%\frac{d {\mathbb Q}}{d {\mathbb P}}^{(\delta)} = 1 + \delta {\Delta} > 0
%\end{eqnarray*}
%for some random variable $\Delta: \Omega \rightarrow {\mathbb R}$ satisfying ${\mathbb E}_{\mathbb P}(\Delta)=0$,
 ${\mathcal A}(\varepsilon;f)\sim O(\varepsilon^\frac{1}{k})$, and  worst-case sensitivity ${\mathcal S}_{\mathbb P}[f(x, \cdot)]$ defined by \eqref{eq:sensitivity-general2}  is a generalized measure of deviation with $g(\varepsilon)=\varepsilon^\frac{1}{k}$.
\end{proposition}
Proposition \ref{prop:deviation} tells us that the expected cost is sensitive to misspecification if its distribution has a large spread. %We now show that different measures of spread correspond to worst-case perturbations under different uncertainty sets.
The spread of the cost distribution is related to robustness because its expected value  is sensitive to small changes in its tail probabilities if it has a large spread, and hence is not robust. However, different uncertainty sets induce different worst-case perturbations and have different measures of spread.  

 
It is well known that for some uncertainty sets ${\mathcal Q}(\varepsilon)$ the ambiguity cost 
\begin{eqnarray}
\label{eq:AC}
{\mathcal A}(\varepsilon; f)  =  \max_{{\mathbb Q}\in{\mathcal Q}(\varepsilon)}{\mathbb E}_{\mathbb Q} \big[f\big] - {\mathbb E}_{\mathbb P}\big[f\big]
\end{eqnarray}
is a generalized measure of deviation and a risk-measure (e.g., Theorem 1 of \cite{rockafellar2006generalized}).  

When the uncertainty set is a constraint on smooth $\phi$-divergence, $g(\varepsilon)=\sqrt{\varepsilon}$ in the definition \eqref{eq:sensitivity-general2} of worst-case sensitivity and \eqref{eq:continuity d} holds with $k=2$ (see \cite{gotoh2026sensitivity} for other examples). 

\section{Penalty version} \label{App:Penalty}
\subsection{Penalty version for smooth $\phi$-divergence}
%\begin{example} \label{example:GKL}
In \cite{gotoh2018robust}, the penalty version of the worst-case problem is used to define worst-case sensitivity. Specifically, a family of worst-case distributions $\{\tilde{\mathsf q}(\delta)\,|\,\delta\geq 0\}$ is given by the solutions of the worst-case problem
\begin{align}
\tilde{\mathsf{q}}(\delta) & := \left\{
\begin{array}{cl}
\begin{displaystyle}\argmax_{\mathsf q}\Big\{  \sum_{i=1}^nq_i f_i -  \frac{1}{\delta} \sum_{i=1}^n {p}_i \phi\Big(\frac{q_i}{{p}_i}\Big)\Big\}, \end{displaystyle}& \delta>0,\\%[10pt]
\mathsf{p}, & \delta=0,
\end{array}\right.
\label{eq:worst-case-Q}
\end{align}
where the parameter $\delta$ determines the penalty on deviations from the nominal. In particular, $\delta=0$ gives the nominal and increasing $\delta$ is analogous to increasing the size of the uncertainty set in \eqref{eq:V}. When the penalty version is used to define the set of worst-case measures, the worst-case expected cost under %${\mathbb Q}(\varepsilon)$
$\tilde{\mathsf{q}}(\delta)$ is linear in the ambiguity parameter
\begin{align*}
%V(\delta) \equiv {\mathbb E}_{{\mathsf q}(\delta)}[f(Y)] = {\mathbb E}_{{\mathsf p}%_n
%}[f(Y)] + \frac{\delta}{\phi''(1)}{\mathbb V}_{{\mathsf p}%_n
%}[f(Y)] + o(\delta)
V(\delta) & \equiv \mathbb{E}_{\tilde{\mathsf{q}}(\delta)}(\mathsf{f}) = \mathbb{E}_{\mathsf{p}}(\mathsf{f})
                +\frac{\delta}{\phi''(1)}\mathbb{V}_{\mathsf{p}}(\mathsf{f}) + o(\delta),
\end{align*}
so the standard definition of sensitivity \eqref{eq:sensitivity-general2} %\eqref{eq:sensitivity-general1} 
can be used:
\begin{align}
%{\mathcal S}_{{\mathsf p}%_n
%}(f) =\lim_{\delta\downarrow 0}\frac{{\mathbb E}_{{\mathsf q}(\delta)}[f(Y)] -{\mathbb E}_{{\mathsf p}%_n
%}[f(Y)]}{\delta} = \frac{1}{\phi''(1)}{\mathbb V}_{{\mathsf p}%_n
%}[f(Y)]
{\mathcal S}_{\mathsf{p}}(\mathsf{f})
 &=\lim_{\delta\downarrow 0}\frac{\mathbb{E}_{\tilde{\mathsf{q}}(\delta)}(\mathsf{f}) -\mathbb{E}_{\mathsf{p}}(\mathsf{f})}{\delta}
  =\frac{1}{\phi''(1)}\mathbb{V}_{\mathsf{p}}(\mathsf{f}).
\label{eq:phi-sensitivity-var}
\end{align}
While this leads to a different sensitivity measure, its qualitative nature is the same as \eqref{eq:sensitivity-general2} with smooth $\phi$-divergence where worst--sensitivity is the standard deviation of the cost distribution.
%\end{example}



\end{document}
